\chapter*{Abstract}
\addcontentsline{toc}{chapter}{\numberline{}Abstract}
A small language model cannot store many facts, so a question about a particular text is better answered by searching a collection first and reading the best passage. This project builds a question-answering system of that kind for Bangla. The reader is BanglaT5, fine-tuned on the BanglaRQA dataset to answer from a passage, or to write nothing when the passage does not contain the answer. The knowledge base is the 3,000 Wikipedia passages of BanglaRQA, searched with BM25, with BGE-M3 dense retrieval, or with a hybrid of the two.

One trained model was tested on the 1,493 held-out questions, first with the correct passage and then with the top passage from each retriever. With the correct passage it reached an exact match of 59.9 and an F1 of 76.1, against 62.4 and 78.1 published for the same model type. It was slightly better than the published model on answerable questions and worse on unanswerable ones. On answerable questions, finding the passage with a retriever cost 4.1 exact-match points with BM25, 6.1 with the hybrid and 10.6 with dense retrieval, and the differences were statistically significant. Almost all of the loss came from questions for which the wrong passage was retrieved. Exact match also undercounts: with the correct passage 82\% of the answers reached an F1 of at least 0.5, while 56\% matched exactly. A small web page shows the system answering questions together with the passage it used.

In a hand-labelled sample of 40 wrong answers given with the correct passage, 78\% differed from the reference only in wording and a further 18\% were partly right; none gave a wrong value.

The study uses general Wikipedia text, one model size, one training run and a library of 3,000 passages, so the results describe this setting and need to be checked on other domains.

\medskip
\noindent\textbf{Keywords:} Bangla, question answering, retrieval-augmented generation, BanglaT5, BM25, dense retrieval, hybrid retrieval, unanswerable questions.
